\begin{papercard}{nielsen1994}
\cf{Stages} \stagetag{W} \stagetag{I}. The book-chapter statement of heuristic evaluation (HE): procedure, the ten revised heuristics, number of evaluators, severity ratings and their reliability, and evaluator expertise. \emph{Added} because Malki and Hutton build on these heuristics. \emph{Note:} page numbers are relative to the appended chapter (printed pp.\ 25--62).

\cf{Definition} A ``discount usability engineering'' method: a small set of evaluators examines the interface and judges compliance with recognised principles (the heuristics) \pc{nielsen1994}{2}{Heuristic evaluation involves having a small set of evaluators}.

\cf{Procedure} Multiple evaluators, because different people find different problems (Figure 2.1: 19 evaluators, 16 problems) \pc{nielsen1994}{2}{different people find different usability problems}. Recommended: three to five evaluators. Each evaluator inspects \emph{alone}, and only after all are done may they communicate and aggregate, to keep evaluations independent \pc{nielsen1994}{2}{Heuristic evaluation is performed by having each individual evaluator inspect}. Recording by written reports or by an observer noting verbalised comments \pc{nielsen1994}{3}{be read and aggregated by an evaluation manager}. Sessions last one to two hours. Unlike user tests, observers may answer domain questions and give hints once the evaluator has noted the problem \pc{nielsen1994}{4}{the evaluators should not be given help}. At least two passes: one for flow and scope, one for specific elements \pc{nielsen1994}{5}{A general recommendation would be that they}. Category-specific heuristics may supplement the general list \pc{nielsen1994}{5}{it is possible to develop category-specific heuristics}.

\cf{The heuristics} Original 1990 list (Table 2.1) \pc{nielsen1994}{5}{Original list of usability heuristics}. Revised set from a factor analysis of 249 problems (Table 2.2): visibility of system status; match with the real world; user control and freedom; consistency and standards; error prevention; recognition rather than recall; flexibility and efficiency; aesthetic and minimalist design; help users recognise, diagnose and recover from errors; help and documentation \pc{nielsen1994}{6}{Revised set of usability heuristics derived from a factor analysis}. For domain-specific systems, give evaluators a usage scenario derived from task analysis \pc{nielsen1994}{7}{One approach that has been applied successfully is}.

\cf{Output} A list of problems, each referenced to the violated heuristic and listed \emph{separately}, with reasons, not mere dislike \pc{nielsen1994}{7}{The output from using the heuristic evaluation method is a list of usability}. HE gives no systematic fixes; an optional debriefing helps redesign \pc{nielsen1994}{8}{One possibility for extending the heuristic evaluation method to provide}.

\cf{Number of evaluators} Single evaluators found 35\,\% of problems on average across six projects \pc{nielsen1994}{8}{single evaluators found only 35 percent of the}. Nielsen--Landauer model: problems found $=N(1-(1-\lambda)^i)$, with $\lambda$ the proportion found by one evaluator (mean 34\,\%) \pc{nielsen1994}{9}{Nielsen and Landauer (1993) present such a model}. A cost-benefit example gives an optimum of four \pc{nielsen1994}{11}{confirming the general observation that heuristic evaluation seems to work best}.

\cf{Case study design} Eleven usability specialists, with domain knowledge checked by a term questionnaire including a \emph{meaningless control term} \pc{nielsen1994}{14}{As a control, the questionnaire also contained an}. Four phases: training, evaluation, debriefing, severity rating \pc{nielsen1994}{14}{The heuristic evaluation consisted of four phases}. Evaluators were not shown screenshots beforehand, to stay unbiased \pc{nielsen1994}{15}{To ensure that the evaluators approached the Integrating System user interface}. One-hour individual sessions, scenario first, then element-level analysis \pc{nielsen1994}{16}{Each evaluator was given one hour to perform the evaluation}. Finding problems and rating them are different processes that should not be interleaved \pc{nielsen1994}{17}{It appears that finding usability problems and the more detailed}.

\cf{Severity ratings and reliability} Collect severity ratings afterwards by questionnaire over the full problem list, independently per evaluator \pc{nielsen1994}{24}{Instead, severity ratings}. Scale 0--4: not a problem, cosmetic, minor, major, catastrophe \pc{nielsen1994}{25}{Five-point rating scale for the severity of usability problems found by}. Inter-rater reliability was \emph{low}: mean pairwise $r=.24$, Kendall's $W=.31$ \pc{nielsen1994}{25}{The average correlation between the severity ratings provided by any two}. Means of several raters are much more reliable (Spearman--Brown), so single-evaluator ratings should not be trusted and 3--4 raters are satisfactory \pc{nielsen1994}{31}{one can definitely conclude that severity}.

\cf{Evaluator expertise} Novices found 22\,\%, usability specialists 41\,\%, and ``double experts'' (usability plus domain) 60\,\% \pc{nielsen1994}{35}{The performance of the}. End users are not recommended as heuristic evaluators.

\cf{Use for this thesis} If heuristics supplement the walkthrough (as in Malki), follow the core protocol: independent evaluation before discussion, problems listed separately with the violated heuristic, two passes, and severity rated afterwards by several evaluators. The low severity reliability is a cautionary data point for any ordinal rating in the thesis, and the ``independent first, discuss later'' rule matches the two-coder requirement.
\end{papercard}
